% Generated draft: requires mathematical and source review.
\documentclass[11pt]{article}
\usepackage[margin=25mm]{geometry}
\usepackage{fontspec}
\setmainfont{Latin Modern Roman}
\usepackage{amsmath,amssymb,mathtools,unicode-math}
\setmathfont{Latin Modern Math}
\usepackage{enumitem,needspace,xurl,hyperref}
\setlength{\parindent}{0pt}
\setlength{\parskip}{5pt}
\setlength{\emergencystretch}{3em}
\newcommand{\N}{\mathbb{N}}
\newcommand{\Z}{\mathbb{Z}}
\newcommand{\Q}{\mathbb{Q}}
\newcommand{\R}{\mathbb{R}}
\newcommand{\C}{\mathbb{C}}
\newcommand{\F}{\mathbb{F}}
\newcommand{\A}{\mathbb{A}}
\newcommand{\PP}{\mathbb{P}}
\newcommand{\E}{\mathbb{E}}
\newcommand{\eps}{\varepsilon}
\newcommand{\dd}{\,\mathrm{d}}
\begin{document}
\section*{UnsolvedMath ID: 30005688}
\textbf{Source problem code:} catalogue entry
\subsection{Definitions and mathematical statement}
The supplied record concerns the continuous Heisenberg group \(G=H_3(\R)\) and a lattice \(\Gamma\leq G\), meaning a discrete cocompact subgroup. In the coordinates used by the source, it specifies a relatively compact left fundamental domain
\[
V=[-1,1)^3
\]
and, for an integer \(\lambda_0\geq 3\), the group automorphism
\[
D_{\lambda_0}(x,y,z)=(\lambda_0x,\lambda_0y,\lambda_0^2z).
\]
The supplied excerpt does not specify the coordinate group law or identify \(\Gamma\) explicitly; these coordinate assertions must therefore be read with the source's unspecified conventions.

Let \(A\) denote the source's alphabet. For a set \(B\), write \(A^B\) for the set of maps from \(B\) to \(A\). The stated substitution rule is a map
\[
S_0:A\longrightarrow A^{D_{\lambda_0}(V)\cap\Gamma}.
\]
The excerpt asserts that this rule extends canonically to coloured patterns and induces a continuous substitution map \(S:A^\Gamma\to A^\Gamma\). For a configuration \(\omega\in A^\Gamma\), its hull is
\[
\Omega_\omega=\overline{\{\gamma\cdot\omega:\gamma\in\Gamma\}},
\]
where the closure is taken in the product topology and \(\gamma\cdot\omega\) denotes the source's translation action. The alphabet assumptions, the precise extension procedure, and the definitions of primitivity and nonperiodicity are not included in the supplied excerpt.

\textbf{Requested determination, with unresolved specifications.} For a nonperiodic configuration \(\omega\) arising from a symbolic substitution system over \(\Gamma\), determine the spectrum
\[
\sigma(H_\omega)
\]
and its spectral distribution function. The original question asks more broadly what can be said about these objects; it does not prescribe a particular form of description.

The supplied record does not define \(H_\omega\), its underlying space, its domain, its coefficients, or its dependence on \(\omega\). It also does not define the spectral distribution function or specify whether this means an integrated density of states, a distribution associated with a particular spectral measure, or another object. Consequently, the displayed request is an incomplete problem specification: a precise operator family, an admissible class of configurations, and a definition and normalization of the distribution function are required before it becomes a determinate mathematical problem. Primitivity occurs in the supplied background assessment, but is not explicitly imposed by the clean statement; its role in the intended question remains unresolved.
\subsection{Short English statement}
The problem asks for a description of spectra and spectral distribution functions of operators associated with nonperiodic Heisenberg substitution configurations. The supplied excerpt describes part of the substitution framework but omits the operator and the definition of its distribution function, so a fully specified mathematical statement requires further source information.
\subsection{Sources}
\begin{itemize}
\item Local UnsolvedMath registry: supplied record with permanent ID \(30005688\), title \emph{Spectra of Operators from Heisenberg Substitutions}, source problem code \texttt{catalogue entry}. The recorded proposed year is 2023; no proposer is supplied.
\item Supplied report citation: \emph{Aspects of Aperiodic Order (2024)}, Oberwolfach Reports, pp.\ 2107--2178, DOI \texttt{10.4171/owr/2023/37}. \url{https://doi.org/10.4171/owr/2023/37}. The background attributes the problem to workshop open-problem discussions documented in Oberwolfach Reports.
\item Supplied literature source, labelled ``Beckus-Hartnick-Pogorzelski'': \url{https://arxiv.org/abs/2109.15210}.
\item Supplied literature source, labelled ``Beckus-Takase'': \url{https://arxiv.org/abs/2110.05763}.
\item Literature assessment checked at \texttt{2026-08-22}: this is supplied dated curation, not independently verified here. It labels the problem partially solved and reports that, for primitive nonperiodic Heisenberg substitutions, the integrated density of states has a Pastur--Shubin limit and periodic spectra approximate the aperiodic spectrum exponentially, while a complete description of the spectrum or distribution function is not known. The record explicitly requests independent verification of these time-sensitive claims and their hypotheses.
\item The registry labels its statement \texttt{corrected\_verified} and says it was reconstructed and checked against the official report. These are supplied source-verification claims, not verification performed for this draft. The original excerpt refers to references [2] and [4] and begins an incomplete theorem attributed to ``Beckus--Hartnick--P.''; its bibliography and complete theorem are not supplied.
\end{itemize}
\end{document}
